\section{Calibration, Simulation Code, and Figures}\label{app:calibration}

This appendix documents the calibration choices, provides key multiplier values
under alternative parameterizations, and describes the simulation code used to
generate all figures in the paper.

\subsection{Baseline calibration}\label{app:baseline-calib}
The baseline calibration is reported in Table~\ref{tab:calibration} in the
main text. The steady-state government spending share $\bar g/\bar y=0.20$
implies $\bar g/\bar c=0.25$, used in $\alpha^{j}=\bar c/(\bar c+\theta^{j}\bar g)$.

\subsection{Key multiplier values under alternative calibrations}\label{app:mult-values}

\begin{table}[ht]
\centering\small
\caption{Key multiplier values across alternative calibrations.}
\label{tab:mult-values}
\begin{threeparttable}
\begin{tabular}{@{}lrrr@{}}
\toprule
Case & {$\chi$} & {$\xi$} & {Multiplier $\mu$} \\
\midrule
Bilbiie separable ($\theta=0$, $\lambda=0.35$) & 2.000 & 0.429 & 1.500 \\
Symmetric complementarity ($\theta=-0.5$, $\lambda=0.35$) & 2.000 & 0.304 & 1.355 \\
Symmetric substitutability ($\theta=+0.3$, $\lambda=0.35$) & 2.000 & 0.498 & 1.581 \\
Heterogeneous internalization ($\theta^{S}=+0.29,\theta^{H}=-0.76,\lambda=0.35$) & 1.874 & 0.530 & 1.539 \\
Estimated upper bound ($\theta^{S}=+0.38,\theta^{H}=-0.90$) & 1.802 & 0.611 & 1.612 \\
Developing economy ($\lambda=0.70,\theta^{H}=-1.20$) & 2.345 & 0.890 & 2.640 \\
Determinacy boundary ($\lambda\chi=1$) & --- & --- & $\infty$ \\
\bottomrule
\end{tabular}
\begin{tablenotes}\footnotesize
\item Note on the determinacy boundary: at $\lambda=0.50$ with $\theta=-0.5$,
$\chi=2.0$ exactly, so $\lambda\chi=1$ and the multiplier diverges. This is
the well-known HANK indeterminacy region identified by \citet{Bilbiie2008}:
when $\lambda\chi\ge 1$, the NK Cross feedback loop is explosive and the
equilibrium is indeterminate. All simulations use $\lambda=0.35$ or
$\lambda=0.70$, both of which are safely below the boundary.
\end{tablenotes}
\end{threeparttable}
\end{table}

\subsection{Figures}\label{app:figures}

All figures in the paper are generated by the Python simulation code in
\path{code/simulation/simulations.py} and
\path{code/hank/hank_jacobian.py}. The code implements all analytical
functions ($\alpha^{j},\Gamma^{j},\chi,\xi,\mu$) directly from the closed-form
expressions in Section~\ref{sec:analytical}, verifies the multiplier ranking
of Corollaries~\ref{cor:sym} and~\ref{cor:het-explicit} numerically, solves
the HANK model via the sequence-space Jacobian method, and generates all
figures referenced in Sections~\ref{sec:quant} and~\ref{sec:policy}.

The full code is open-source and available in the replication repository
(see README.md). Dependencies: \texttt{numpy}, \texttt{matplotlib},
\texttt{scipy}, \texttt{pandas}.

\subsection{Simulation code}\label{app:sim-code}

The simulation is organized into four modules:
\begin{itemize}[leftmargin=1.6em]
  \item \path{code/simulation/analytical.py}: computes $\chi$, $\xi$, $\mu$
    from the closed-form expressions.
  \item \path{code/simulation/figures.py}: generates all paper figures.
  \item \path{code/hank/hank_jacobian.py}: solves the HANK model via the
    sequence-space Jacobian method.
  \item \path{code/empirical/cex_analysis.do}: implements the CEX
    diff-in-diff estimation of Section~\ref{sec:empirical}.
\end{itemize}

Run \texttt{python3 run\_all.py} from the repository root to reproduce all
quantitative results and figures in the paper.
